% =================================
% NGÂN HÀNG CÂU HỎI MẪU
% Chỉ chứa các lệnh \caugiai{...}{...}
% =================================

\caugiai
{
Đồ thị của hàm số nào dưới đây có dạng đường cong như hình vẽ bên?
(Giả sử có hình vẽ đồ thị bậc 3)
\bonpa
{\(y=x^3-3x\)}
{\(y=-x^3+3x\)}
{\(y=x^4-2x^2\)}
{\(y=-x^4+2x^2\)}
}
{
Đường cong có dạng của đồ thị hàm số bậc ba với hệ số \(a>0\). 

Vậy hàm số đó là \(y=x^3-3x\). Chọn A.
}

\caugiai
{
Cho hàm số \(y=f(x)\) có bảng biến thiên như sau. Hàm số đã cho đồng biến trên khoảng nào dưới đây?
\bonpa
{\((0;+\infty)\)}
{\((-\infty;-2)\)}
{\((-2;0)\)}
{\((0;2)\)}
}
{
Dựa vào bảng biến thiên (giả định), ta thấy hàm số đồng biến trên khoảng \((-2;0)\). Chọn C.
}

\caugiai
{
Giá trị lớn nhất của hàm số \(f(x) = x^3 - 3x + 2\) trên đoạn \([-1; 2]\) bằng bao nhiêu?
\bonpa
{\(4\)}
{\(2\)}
{\(0\)}
{\(-2\)}
}
{
Ta có \(f'(x) = 3x^2 - 3\). Cho \(f'(x) = 0 \Leftrightarrow x = \pm 1\).
Trên đoạn \([-1; 2]\), ta tính:
\(f(-1) = 4\)
\(f(1) = 0\)
\(f(2) = 4\)

Vậy giá trị lớn nhất bằng 4. Chọn A.
}

\caugiai
{
Đường tiệm cận đứng của đồ thị hàm số \(y=\dfrac{2x-1}{x-1}\) là đường thẳng có phương trình:
\bonpa
{\(x=1\)}
{\(x=2\)}
{\(y=1\)}
{\(y=2\)}
}
{
Cho mẫu số bằng 0 ta được \(x-1=0 \Leftrightarrow x=1\). 

Vì nghiệm tử khác 1 nên \(x=1\) là đường tiệm cận đứng. Chọn A.
}

\caugiai
{
Tìm tất cả các giá trị thực của tham số \(m\) để hàm số \(y=x^3-3x^2+mx+1\) đồng biến trên \(\mathbb{R}\).
\bonpa
{\(m \ge 3\)}
{\(m > 3\)}
{\(m \le 3\)}
{\(m < 3\)}
}
{
Ta có \(y'=3x^2-6x+m\).
Để hàm số đồng biến trên \(\mathbb{R}\) thì \(y' \ge 0, \forall x \in \mathbb{R}\).
Điều này tương đương với \(\Delta' = 9 - 3m \le 0 \Leftrightarrow m \ge 3\). Chọn A.
}

\caugiai
{
Đạo hàm của hàm số \(y = x^4 - 2x^2 + 1\) là:
\bonpa
{\(y' = 4x^3 - 4x\)}
{\(y' = 4x^3 - 2x\)}
{\(y' = x^3 - 4x\)}
{\(y' = 4x^3 + 4x\)}
}
{
Ta có \(y' = 4x^3 - 4x\). Chọn A.
}

\caugiai
{
Tập xác định của hàm số \(y = \log_2(x-1)\) là:
\bonpa
{\((1; +\infty)\)}
{\([1; +\infty)\)}
{\((-\infty; 1)\)}
{\(\mathbb{R} \setminus \{1\}\)}
}
{
Điều kiện xác định: \(x-1 > 0 \Leftrightarrow x > 1\). 
Vậy tập xác định là \(D = (1; +\infty)\). Chọn A.
}

\caugiai
{
Thể tích \(V\) của khối chóp có diện tích đáy \(B=6\) và chiều cao \(h=4\) là:
\bonpa
{\(V = 8\)}
{\(V = 24\)}
{\(V = 12\)}
{\(V = 4\)}
}
{
Thể tích khối chóp là \(V = \dfrac{1}{3}Bh = \dfrac{1}{3} \cdot 6 \cdot 4 = 8\). Chọn A.
}

\caugiai
{
Cho cấp số cộng \((u_n)\) có \(u_1 = 2\) và công sai \(d = 3\). Giá trị của \(u_3\) bằng:
\bonpa
{\(8\)}
{\(5\)}
{\(11\)}
{\(6\)}
}
{
Ta có \(u_3 = u_1 + 2d = 2 + 2 \cdot 3 = 8\). Chọn A.
}

\caugiai
{
Tích phân \(\displaystyle\int_0^1 2x \dd x\) bằng:
\bonpa
{\(1\)}
{\(2\)}
{\(0\)}
{\(\dfrac{1}{2}\)}
}
{
Ta có \(\displaystyle\int_0^1 2x \dd x = x^2 \Big|_0^1 = 1^2 - 0^2 = 1\). Chọn A.
}

\caugiai
{
Số phức liên hợp của số phức \(z = 3 - 4i\) là:
\bonpa
{\(\bar{z} = 3 + 4i\)}
{\(\bar{z} = -3 - 4i\)}
{\(\bar{z} = -3 + 4i\)}
{\(\bar{z} = 4 - 3i\)}
}
{
Số phức liên hợp của \(z = a + bi\) là \(\bar{z} = a - bi\). 
Do đó số phức liên hợp của \(z = 3 - 4i\) là \(\bar{z} = 3 + 4i\). Chọn A.
}

\caugiai
{
Trong không gian \(Oxyz\), mặt cầu \((S): (x-1)^2 + (y+2)^2 + (z-3)^2 = 9\) có bán kính \(R\) bằng:
\bonpa
{\(3\)}
{\(9\)}
{\(81\)}
{\(6\)}
}
{
Phương trình mặt cầu có dạng \((x-a)^2+(y-b)^2+(z-c)^2=R^2\). 
Ta có \(R^2 = 9 \Rightarrow R = 3\). Chọn A.
}

\caugiai
{
Cho hàm số \(f(x)\) liên tục trên \(\mathbb{R}\) thỏa mãn \(\displaystyle\int_1^2 f(x) \dd x = 3\) và \(\displaystyle\int_1^3 f(x) \dd x = 5\). Tích phân \(\displaystyle\int_2^3 f(x) \dd x\) bằng:
\bonpa
{\(2\)}
{\(8\)}
{\(-2\)}
{\(15\)}
}
{
Ta có \(\displaystyle\int_1^3 f(x) \dd x = \int_1^2 f(x) \dd x + \int_2^3 f(x) \dd x\).
Suy ra \(\displaystyle\int_2^3 f(x) \dd x = 5 - 3 = 2\). Chọn A.
}

\caugiai
{
Nghiệm của phương trình \(2^{x-1} = 8\) là:
\bonpa
{\(x = 4\)}
{\(x = 3\)}
{\(x = 2\)}
{\(x = 5\)}
}
{
Ta có \(2^{x-1} = 8 \Leftrightarrow 2^{x-1} = 2^3 \Leftrightarrow x - 1 = 3 \Leftrightarrow x = 4\). Chọn A.
}

\caugiai
{
Hoành độ giao điểm của đồ thị hàm số \(y=x^2-2x\) và đường thẳng \(y=x-2\) là:
\bonpa
{\(x=1, x=2\)}
{\(x=-1, x=2\)}
{\(x=1, x=-2\)}
{\(x=-1, x=-2\)}
}
{
Phương trình hoành độ giao điểm:
\(x^2 - 2x = x - 2 \Leftrightarrow x^2 - 3x + 2 = 0 \Leftrightarrow \left[ \begin{array}{l} x = 1 \\ x = 2 \end{array} \right.\)
Vậy có hai hoành độ giao điểm là \(x=1\) và \(x=2\). Chọn A.
}
